
Researchers spend 2-4 weeks manually reviewing benchmark papers to determine suitability for hypothesis validation. This work investigates whether benchmark design features create systematic constraints that enable predictive coverage analysis. Using a pilot study of 20 diverse benchmarks (vision, language, audio, multimodal), design features (task formulation, evaluation metrics, data modality, dataset characteristics) are extracted via a standardized protocol achieving Cohen's kappa $\geq 0.917$. SentenceBERT embeddings of feature descriptions are clustered using k-means, discovering four coverage families organized primarily by data modality rather than task type. Citation co-occurrence within coverage families reaches 78.07\% (Jaccard similarity), significantly exceeding random baseline (19.61\%, p<0.001). Historical validation (pre-2023 benchmarks predicting 2023-2024 patterns) demonstrates temporal persistence over 1-2 year publication cycles. These findings suggest that modality-driven metric constraints create coverage patterns predictive of benchmark adoption, with potential to reduce manual selection effort.
